\section{Bewertungsgrundlagen und Metriken}
Die Zuverlässigkeit, die gemessen werden soll, kann in vier Oberkategorien aufgeteilt werden.

Die erste Kategorie ist die Verfügbarkeit. Dabei wird gemessen, wie gut der Service erreichbar ist. Das wird vor allem daran gemessen, ob der Service bei Ausfällen Request bearbeiten kann.
Es wird auch gemessen, wie lange ein Service bei einem Ausfall offline bleibt.
 
Die zweite Kategorie ist die Fehlertoleranz. Dabei soll gemessen werden, wie mit Fehlern bei Ausfällen und bei Last umgegangen wird.
Dabei wird vor allem die Error Rate gemessen, die während der Szenarien auftritt, sowie die Frage, ob es zu einem Neustart kommt.

Die dritte Kategorie ist die Stabilität. Dabei wird das Verhalten unter steigender Last getestet.
Dafür wird zueinen der CPU- und RAM-Verbrauch der Anwendung, aber auch vom gesamten Systems erfasst. Damit soll auch getestet werden, wie die Umgebung die Last skaliert.

Die vierte Kategorie ist die Wiederherstellbarkeit, also die Zeit, die benötigt wird, um den Service nach einem Ausfall wiederherzustellen.
Dabei wird die TTR (Time to Recovery) gemessen, also die Zeit, die ein System zum Wiederherstellen des Services braucht.

Eine Kategorie, die nur teilweise betrachtet wird, ist die Performance. Zwar beeinflussen sich Performance und Zuverlässigkeit gegenseitig, dennoch sind das zwei verschiedene Aspekte, die meist Hand in Hand gehen.
Dennoch werden die p95-Latenz und die durchschnittliche Response Time mitgemessen.

Mithilfe eines eigens erstellten Scoring-Verfahrens soll die Zuverlässigkeit bei jedem Test von 0 bis 100 bewertet werden.
Das Scoring beginnt mit einem Wert von 100 und nimmt ab, wenn die aufgenommenen Metriken auf eine eher schlechte Zuverlässigkeit hinweisen. Die Untergrenze bei jedem Scoring ist 0.
Dieses Scoring soll dabei helfen, die beiden Systeme besser vergleichen zu können. Je nachdem, was getestet wird, kann sich das Scoring unterscheiden.

Die Metriken sind normalisiert, damit beide Systeme besser verglichen werden können.